
% Please see Figure~\ref{fig:quality_filter} for the comparison of video quality before and after this filtering process.
\begin{figure*}[!t]
    \centering
    \includegraphics[width=1.0\linewidth]{figs/supp/quality_filter.png}
    \caption{Videos we keep and deleted based on different quality filters}
    \label{fig:quality_filter}
\end{figure*}

%%%%%%%%% BODY TEXT
\section{Dataset}
\subsection{Video Quality Filter Details}
This study incorporates a comprehensive video quality filtering process grounded in several critical metrics: luminance, blurriness, aesthetic appeal, global and local motion, and overall technical quality. Each metric is evaluated against specific thresholds to retain only high-quality video content for further analysis.

\noindent\textbf{(1) Luminance.}  
Luminance is calculated using the formula: \(0.2126 \times R + 0.7152 \times G + 0.0722 \times B\), where \(R\), \(G\), and \(B\) denote the pixel values of the red, green, and blue channels, respectively. Videos are retained if their luminance values fall within the range of [10, 210].

\noindent\textbf{(2) Blur.}  
The level of blur in the video is assessed via the Cumulative Probability of Blur Detection algorithm \cite{narvekar2011no}, which analyzes edge feature distributions. Frames are extracted from the video, converted to grayscale, and the average, maximum, and minimum blur values are computed. Videos are filtered based on blur values within the range of [0.05, 1.0).

\noindent\textbf{(3) Aesthetic Quality.}  
Aesthetic appeal is measured using the CLIP-based Aesthetics Predictor \cite{aesthetic_predictor}. This method excludes videos with aesthetic scores below 4.8, thereby ensuring a satisfactory visual experience.

\noindent\textbf{(4) Global and Local Motion.}  
Motion quality is evaluated using optical flow analysis, specifically through the UniMatch algorithm \cite{unimatch}, which provides scores for both local and global motion. The filtering criteria for these scores are set within the range of [0.5, 20].

\noindent\textbf{(5) Technical Quality.}  
Lastly, the overall technical quality of the video is assessed using the DOVER model \cite{wu2023exploring}, filtering out videos with scores below 0.09.

\subsection{Human Motion Generation Details}

\noindent\textbf{Textual Condition.}  
Figure~\ref{fig:supp_caption} illustrates examples of various caption formats. We provide final captions in short, long, and structured formats to capture different contextual details.

\noindent\textbf{Skeleton Sequence.}  
Figure~\ref{fig:supp_skeleton} presents examples of skeleton annotations extracted from videos using DWpose \cite{yang2023effective}. These annotations offer a structured representation of joint positions that inform the movement of the generated figures.

\noindent\textbf{Speech Audio.}  
Figure~\ref{fig:supp_speech} shows examples of speech audio from the videos. We utilize SyncNet \cite{raina2022syncnet} to evaluate the synchronization between the subject's speech audio and their corresponding lip movements.


\begin{figure*}[!htbp]
    \centering
    \includegraphics[width=1.0\linewidth]{figs/supp/supp_caption.pdf} 
    \caption{The illustration of textual captions and corresponding types (long, short, and structured captions). }
    \label{fig:supp_caption} 
\end{figure*}


\begin{figure*}[!htbp]
    \centering
    \includegraphics[width=1.0\linewidth]{figs/supp/supp_skeleton.pdf} 
    \caption{The illustration of human skeleton sequences with respect to given videos.}
    \label{fig:supp_skeleton} 
\end{figure*}

\begin{figure*}[!htbp]
    \centering
    \includegraphics[width=1.0\linewidth]{figs/supp/supp_speech.pdf} 
    \caption{The illustration of speech audio with respect to given videos. Left: video screenshot; right: speech script.}
    \label{fig:supp_speech} 
\end{figure*}


\section{Experiments}

\subsection{Training on OpenHumanVid}
We utilize a large-scale, high-quality human video dataset, implementing the training data strategies outlined in the paper. 
These strategies include a higher video sampling rate, human appearance text-video filter, human motion text-video filter, and face motion text-video filter. 
As a result, the extended model demonstrates a significant improvement over the baseline following further pretraining.

% \begin{figure*}
%     \centering
%     \includegraphics[width=1\linewidth]{figs/consistency2.pdf}
%     \caption{Face and body consistency comparison between baseline and ours. }
%     \label{fig:consistency}
% \end{figure*}

% \begin{figure*}
%     \centering
%     \includegraphics[width=1\linewidth]{figs/semantics2.pdf}
%     \caption{Face and body semantics comparison between baseline and ours.}
%     \label{fig:semantics}
% \end{figure*}


\subsection{Limitations and Future Works}
Despite the advancements presented by OpenHumanVid, several limitations warrant consideration. 
Firstly, the reliance on existing multimodal models for caption generation may restrict the diversity and richness of the descriptions, potentially impacting the comprehensiveness of video-text alignment. 
Additionally, while the dataset incorporates various human motion conditions, the inherent variability of human actions and expressions across different cultures and contexts may not be fully captured. 
Future work could focus on enhancing the caption generation process through the development of more sophisticated models with improved instruction-following capabilities. 
Furthermore, expanding the dataset to include a wider array of human identities and motion scenarios, along with more granular textual prompts, could facilitate the generation of videos that better reflect the complexities of human behavior and interaction in diverse contexts.

\section{Safety Considerations}
The development of OpenHumanVid introduces several social risks, particularly concerning privacy, representation, and the potential for misuse. The dataset's inclusion of diverse human identities and movements necessitates rigorous ethical considerations to prevent the reinforcement of stereotypes or biases in generated videos. Additionally, there is a risk of inappropriate usage in contexts such as deepfakes or unauthorized portrayals of individuals. To mitigate these issues, we implement strict data governance protocols, ensuring that all content adheres to ethical standards and is used responsibly. Furthermore, we advocate for transparency in the dataset's application, emphasizing the necessity of guidelines for researchers and practitioners to promote responsible use while safeguarding individual rights and societal norms.

%%%%%%%%% REFERENCES
% {\small
% \bibliographystyle{ieee_fullname}
% \bibliography{egbib}
% }
% 
% Please see Figure~\ref{fig:quality_filter} for the comparison of video quality before and after this filtering process.
\begin{figure*}[!t]
    \centering
    \includegraphics[width=1.0\linewidth]{figs/supp/quality_filter.png}
    \caption{Videos we keep and deleted based on different quality filters}
    \label{fig:quality_filter}
\end{figure*}

%%%%%%%%% BODY TEXT
\section{Dataset}
\subsection{Video Quality Filter Details}
This study incorporates a comprehensive video quality filtering process grounded in several critical metrics: luminance, blurriness, aesthetic appeal, global and local motion, and overall technical quality. Each metric is evaluated against specific thresholds to retain only high-quality video content for further analysis.

\noindent\textbf{(1) Luminance.}  
Luminance is calculated using the formula: \(0.2126 \times R + 0.7152 \times G + 0.0722 \times B\), where \(R\), \(G\), and \(B\) denote the pixel values of the red, green, and blue channels, respectively. Videos are retained if their luminance values fall within the range of [10, 210].

\noindent\textbf{(2) Blur.}  
The level of blur in the video is assessed via the Cumulative Probability of Blur Detection algorithm \cite{narvekar2011no}, which analyzes edge feature distributions. Frames are extracted from the video, converted to grayscale, and the average, maximum, and minimum blur values are computed. Videos are filtered based on blur values within the range of [0.05, 1.0).

\noindent\textbf{(3) Aesthetic Quality.}  
Aesthetic appeal is measured using the CLIP-based Aesthetics Predictor \cite{aesthetic_predictor}. This method excludes videos with aesthetic scores below 4.8, thereby ensuring a satisfactory visual experience.

\noindent\textbf{(4) Global and Local Motion.}  
Motion quality is evaluated using optical flow analysis, specifically through the UniMatch algorithm \cite{unimatch}, which provides scores for both local and global motion. The filtering criteria for these scores are set within the range of [0.5, 20].

\noindent\textbf{(5) Technical Quality.}  
Lastly, the overall technical quality of the video is assessed using the DOVER model \cite{wu2023exploring}, filtering out videos with scores below 0.09.

\subsection{Human Motion Generation Details}

\noindent\textbf{Textual Condition.}  
Figure~\ref{fig:supp_caption} illustrates examples of various caption formats. We provide final captions in short, long, and structured formats to capture different contextual details.

\noindent\textbf{Skeleton Sequence.}  
Figure~\ref{fig:supp_skeleton} presents examples of skeleton annotations extracted from videos using DWpose \cite{yang2023effective}. These annotations offer a structured representation of joint positions that inform the movement of the generated figures.

\noindent\textbf{Speech Audio.}  
Figure~\ref{fig:supp_speech} shows examples of speech audio from the videos. We utilize SyncNet \cite{raina2022syncnet} to evaluate the synchronization between the subject's speech audio and their corresponding lip movements.


\begin{figure*}[!htbp]
    \centering
    \includegraphics[width=1.0\linewidth]{figs/supp/supp_caption.pdf} 
    \caption{The illustration of textual captions and corresponding types (long, short, and structured captions). }
    \label{fig:supp_caption} 
\end{figure*}


\begin{figure*}[!htbp]
    \centering
    \includegraphics[width=1.0\linewidth]{figs/supp/supp_skeleton.pdf} 
    \caption{The illustration of human skeleton sequences with respect to given videos.}
    \label{fig:supp_skeleton} 
\end{figure*}

\begin{figure*}[!htbp]
    \centering
    \includegraphics[width=1.0\linewidth]{figs/supp/supp_speech.pdf} 
    \caption{The illustration of speech audio with respect to given videos. Left: video screenshot; right: speech script.}
    \label{fig:supp_speech} 
\end{figure*}


\section{Experiments}

\subsection{Training on OpenHumanVid}
We utilize a large-scale, high-quality human video dataset, implementing the training data strategies outlined in the paper. 
These strategies include a higher video sampling rate, human appearance text-video filter, human motion text-video filter, and face motion text-video filter. 
As a result, the extended model demonstrates a significant improvement over the baseline following further pretraining.

% \begin{figure*}
%     \centering
%     \includegraphics[width=1\linewidth]{figs/consistency2.pdf}
%     \caption{Face and body consistency comparison between baseline and ours. }
%     \label{fig:consistency}
% \end{figure*}

% \begin{figure*}
%     \centering
%     \includegraphics[width=1\linewidth]{figs/semantics2.pdf}
%     \caption{Face and body semantics comparison between baseline and ours.}
%     \label{fig:semantics}
% \end{figure*}


\subsection{Limitations and Future Works}
Despite the advancements presented by OpenHumanVid, several limitations warrant consideration. 
Firstly, the reliance on existing multimodal models for caption generation may restrict the diversity and richness of the descriptions, potentially impacting the comprehensiveness of video-text alignment. 
Additionally, while the dataset incorporates various human motion conditions, the inherent variability of human actions and expressions across different cultures and contexts may not be fully captured. 
Future work could focus on enhancing the caption generation process through the development of more sophisticated models with improved instruction-following capabilities. 
Furthermore, expanding the dataset to include a wider array of human identities and motion scenarios, along with more granular textual prompts, could facilitate the generation of videos that better reflect the complexities of human behavior and interaction in diverse contexts.

\section{Safety Considerations}
The development of OpenHumanVid introduces several social risks, particularly concerning privacy, representation, and the potential for misuse. The dataset's inclusion of diverse human identities and movements necessitates rigorous ethical considerations to prevent the reinforcement of stereotypes or biases in generated videos. Additionally, there is a risk of inappropriate usage in contexts such as deepfakes or unauthorized portrayals of individuals. To mitigate these issues, we implement strict data governance protocols, ensuring that all content adheres to ethical standards and is used responsibly. Furthermore, we advocate for transparency in the dataset's application, emphasizing the necessity of guidelines for researchers and practitioners to promote responsible use while safeguarding individual rights and societal norms.

%%%%%%%%% REFERENCES
% {\small
% \bibliographystyle{ieee_fullname}
% \bibliography{egbib}
% }
% 
% Please see Figure~\ref{fig:quality_filter} for the comparison of video quality before and after this filtering process.
\begin{figure*}[!t]
    \centering
    \includegraphics[width=1.0\linewidth]{figs/supp/quality_filter.png}
    \caption{Videos we keep and deleted based on different quality filters}
    \label{fig:quality_filter}
\end{figure*}

%%%%%%%%% BODY TEXT
\section{Dataset}
\subsection{Video Quality Filter Details}
This study incorporates a comprehensive video quality filtering process grounded in several critical metrics: luminance, blurriness, aesthetic appeal, global and local motion, and overall technical quality. Each metric is evaluated against specific thresholds to retain only high-quality video content for further analysis.

\noindent\textbf{(1) Luminance.}  
Luminance is calculated using the formula: \(0.2126 \times R + 0.7152 \times G + 0.0722 \times B\), where \(R\), \(G\), and \(B\) denote the pixel values of the red, green, and blue channels, respectively. Videos are retained if their luminance values fall within the range of [10, 210].

\noindent\textbf{(2) Blur.}  
The level of blur in the video is assessed via the Cumulative Probability of Blur Detection algorithm \cite{narvekar2011no}, which analyzes edge feature distributions. Frames are extracted from the video, converted to grayscale, and the average, maximum, and minimum blur values are computed. Videos are filtered based on blur values within the range of [0.05, 1.0).

\noindent\textbf{(3) Aesthetic Quality.}  
Aesthetic appeal is measured using the CLIP-based Aesthetics Predictor \cite{aesthetic_predictor}. This method excludes videos with aesthetic scores below 4.8, thereby ensuring a satisfactory visual experience.

\noindent\textbf{(4) Global and Local Motion.}  
Motion quality is evaluated using optical flow analysis, specifically through the UniMatch algorithm \cite{unimatch}, which provides scores for both local and global motion. The filtering criteria for these scores are set within the range of [0.5, 20].

\noindent\textbf{(5) Technical Quality.}  
Lastly, the overall technical quality of the video is assessed using the DOVER model \cite{wu2023exploring}, filtering out videos with scores below 0.09.

\subsection{Human Motion Generation Details}

\noindent\textbf{Textual Condition.}  
Figure~\ref{fig:supp_caption} illustrates examples of various caption formats. We provide final captions in short, long, and structured formats to capture different contextual details.

\noindent\textbf{Skeleton Sequence.}  
Figure~\ref{fig:supp_skeleton} presents examples of skeleton annotations extracted from videos using DWpose \cite{yang2023effective}. These annotations offer a structured representation of joint positions that inform the movement of the generated figures.

\noindent\textbf{Speech Audio.}  
Figure~\ref{fig:supp_speech} shows examples of speech audio from the videos. We utilize SyncNet \cite{raina2022syncnet} to evaluate the synchronization between the subject's speech audio and their corresponding lip movements.


\begin{figure*}[!htbp]
    \centering
    \includegraphics[width=1.0\linewidth]{figs/supp/supp_caption.pdf} 
    \caption{The illustration of textual captions and corresponding types (long, short, and structured captions). }
    \label{fig:supp_caption} 
\end{figure*}


\begin{figure*}[!htbp]
    \centering
    \includegraphics[width=1.0\linewidth]{figs/supp/supp_skeleton.pdf} 
    \caption{The illustration of human skeleton sequences with respect to given videos.}
    \label{fig:supp_skeleton} 
\end{figure*}

\begin{figure*}[!htbp]
    \centering
    \includegraphics[width=1.0\linewidth]{figs/supp/supp_speech.pdf} 
    \caption{The illustration of speech audio with respect to given videos. Left: video screenshot; right: speech script.}
    \label{fig:supp_speech} 
\end{figure*}


\section{Experiments}

\subsection{Training on OpenHumanVid}
We utilize a large-scale, high-quality human video dataset, implementing the training data strategies outlined in the paper. 
These strategies include a higher video sampling rate, human appearance text-video filter, human motion text-video filter, and face motion text-video filter. 
As a result, the extended model demonstrates a significant improvement over the baseline following further pretraining.

% \begin{figure*}
%     \centering
%     \includegraphics[width=1\linewidth]{figs/consistency2.pdf}
%     \caption{Face and body consistency comparison between baseline and ours. }
%     \label{fig:consistency}
% \end{figure*}

% \begin{figure*}
%     \centering
%     \includegraphics[width=1\linewidth]{figs/semantics2.pdf}
%     \caption{Face and body semantics comparison between baseline and ours.}
%     \label{fig:semantics}
% \end{figure*}


\subsection{Limitations and Future Works}
Despite the advancements presented by OpenHumanVid, several limitations warrant consideration. 
Firstly, the reliance on existing multimodal models for caption generation may restrict the diversity and richness of the descriptions, potentially impacting the comprehensiveness of video-text alignment. 
Additionally, while the dataset incorporates various human motion conditions, the inherent variability of human actions and expressions across different cultures and contexts may not be fully captured. 
Future work could focus on enhancing the caption generation process through the development of more sophisticated models with improved instruction-following capabilities. 
Furthermore, expanding the dataset to include a wider array of human identities and motion scenarios, along with more granular textual prompts, could facilitate the generation of videos that better reflect the complexities of human behavior and interaction in diverse contexts.

\section{Safety Considerations}
The development of OpenHumanVid introduces several social risks, particularly concerning privacy, representation, and the potential for misuse. The dataset's inclusion of diverse human identities and movements necessitates rigorous ethical considerations to prevent the reinforcement of stereotypes or biases in generated videos. Additionally, there is a risk of inappropriate usage in contexts such as deepfakes or unauthorized portrayals of individuals. To mitigate these issues, we implement strict data governance protocols, ensuring that all content adheres to ethical standards and is used responsibly. Furthermore, we advocate for transparency in the dataset's application, emphasizing the necessity of guidelines for researchers and practitioners to promote responsible use while safeguarding individual rights and societal norms.

%%%%%%%%% REFERENCES
% {\small
% \bibliographystyle{ieee_fullname}
% \bibliography{egbib}
% }
% 
% Please see Figure~\ref{fig:quality_filter} for the comparison of video quality before and after this filtering process.
\begin{figure*}[!t]
    \centering
    \includegraphics[width=1.0\linewidth]{figs/supp/quality_filter.png}
    \caption{Videos we keep and deleted based on different quality filters}
    \label{fig:quality_filter}
\end{figure*}

%%%%%%%%% BODY TEXT
\section{Dataset}
\subsection{Video Quality Filter Details}
This study incorporates a comprehensive video quality filtering process grounded in several critical metrics: luminance, blurriness, aesthetic appeal, global and local motion, and overall technical quality. Each metric is evaluated against specific thresholds to retain only high-quality video content for further analysis.

\noindent\textbf{(1) Luminance.}  
Luminance is calculated using the formula: \(0.2126 \times R + 0.7152 \times G + 0.0722 \times B\), where \(R\), \(G\), and \(B\) denote the pixel values of the red, green, and blue channels, respectively. Videos are retained if their luminance values fall within the range of [10, 210].

\noindent\textbf{(2) Blur.}  
The level of blur in the video is assessed via the Cumulative Probability of Blur Detection algorithm \cite{narvekar2011no}, which analyzes edge feature distributions. Frames are extracted from the video, converted to grayscale, and the average, maximum, and minimum blur values are computed. Videos are filtered based on blur values within the range of [0.05, 1.0).

\noindent\textbf{(3) Aesthetic Quality.}  
Aesthetic appeal is measured using the CLIP-based Aesthetics Predictor \cite{aesthetic_predictor}. This method excludes videos with aesthetic scores below 4.8, thereby ensuring a satisfactory visual experience.

\noindent\textbf{(4) Global and Local Motion.}  
Motion quality is evaluated using optical flow analysis, specifically through the UniMatch algorithm \cite{unimatch}, which provides scores for both local and global motion. The filtering criteria for these scores are set within the range of [0.5, 20].

\noindent\textbf{(5) Technical Quality.}  
Lastly, the overall technical quality of the video is assessed using the DOVER model \cite{wu2023exploring}, filtering out videos with scores below 0.09.

\subsection{Human Motion Generation Details}

\noindent\textbf{Textual Condition.}  
Figure~\ref{fig:supp_caption} illustrates examples of various caption formats. We provide final captions in short, long, and structured formats to capture different contextual details.

\noindent\textbf{Skeleton Sequence.}  
Figure~\ref{fig:supp_skeleton} presents examples of skeleton annotations extracted from videos using DWpose \cite{yang2023effective}. These annotations offer a structured representation of joint positions that inform the movement of the generated figures.

\noindent\textbf{Speech Audio.}  
Figure~\ref{fig:supp_speech} shows examples of speech audio from the videos. We utilize SyncNet \cite{raina2022syncnet} to evaluate the synchronization between the subject's speech audio and their corresponding lip movements.


\begin{figure*}[!htbp]
    \centering
    \includegraphics[width=1.0\linewidth]{figs/supp/supp_caption.pdf} 
    \caption{The illustration of textual captions and corresponding types (long, short, and structured captions). }
    \label{fig:supp_caption} 
\end{figure*}


\begin{figure*}[!htbp]
    \centering
    \includegraphics[width=1.0\linewidth]{figs/supp/supp_skeleton.pdf} 
    \caption{The illustration of human skeleton sequences with respect to given videos.}
    \label{fig:supp_skeleton} 
\end{figure*}

\begin{figure*}[!htbp]
    \centering
    \includegraphics[width=1.0\linewidth]{figs/supp/supp_speech.pdf} 
    \caption{The illustration of speech audio with respect to given videos. Left: video screenshot; right: speech script.}
    \label{fig:supp_speech} 
\end{figure*}


\section{Experiments}

\subsection{Training on OpenHumanVid}
We utilize a large-scale, high-quality human video dataset, implementing the training data strategies outlined in the paper. 
These strategies include a higher video sampling rate, human appearance text-video filter, human motion text-video filter, and face motion text-video filter. 
As a result, the extended model demonstrates a significant improvement over the baseline following further pretraining.

% \begin{figure*}
%     \centering
%     \includegraphics[width=1\linewidth]{figs/consistency2.pdf}
%     \caption{Face and body consistency comparison between baseline and ours. }
%     \label{fig:consistency}
% \end{figure*}

% \begin{figure*}
%     \centering
%     \includegraphics[width=1\linewidth]{figs/semantics2.pdf}
%     \caption{Face and body semantics comparison between baseline and ours.}
%     \label{fig:semantics}
% \end{figure*}


\subsection{Limitations and Future Works}
Despite the advancements presented by OpenHumanVid, several limitations warrant consideration. 
Firstly, the reliance on existing multimodal models for caption generation may restrict the diversity and richness of the descriptions, potentially impacting the comprehensiveness of video-text alignment. 
Additionally, while the dataset incorporates various human motion conditions, the inherent variability of human actions and expressions across different cultures and contexts may not be fully captured. 
Future work could focus on enhancing the caption generation process through the development of more sophisticated models with improved instruction-following capabilities. 
Furthermore, expanding the dataset to include a wider array of human identities and motion scenarios, along with more granular textual prompts, could facilitate the generation of videos that better reflect the complexities of human behavior and interaction in diverse contexts.

\section{Safety Considerations}
The development of OpenHumanVid introduces several social risks, particularly concerning privacy, representation, and the potential for misuse. The dataset's inclusion of diverse human identities and movements necessitates rigorous ethical considerations to prevent the reinforcement of stereotypes or biases in generated videos. Additionally, there is a risk of inappropriate usage in contexts such as deepfakes or unauthorized portrayals of individuals. To mitigate these issues, we implement strict data governance protocols, ensuring that all content adheres to ethical standards and is used responsibly. Furthermore, we advocate for transparency in the dataset's application, emphasizing the necessity of guidelines for researchers and practitioners to promote responsible use while safeguarding individual rights and societal norms.

%%%%%%%%% REFERENCES
% {\small
% \bibliographystyle{ieee_fullname}
% \bibliography{egbib}
% }
% \input{sections/appendix}

\end{document}

\end{document}

\end{document}

\end{document}